\section{Preliminary DEM and graph-based analysis}
\label{sec:preliminary_dem_gnn}

\subsection{Agglomerate generation and numerical design}

The preliminary data set consists of 135 monodisperse agglomerates generated
with the local-filling-factor algorithm of Ringl and Urbassek. In this method,
a new primary particle is attached preferentially to a particle with a low
local filling factor. The detector radius is written as
\begin{equation}
    r_{\sigma}=m_{\sigma}r_p,
\end{equation}
where $r_p$ is the primary-particle radius and $m_{\sigma}$ is the detector-radius
multiplier. The latter controls the generated morphology but is not itself the
fractal dimension. Consequently, the requested values of $D_f$ should be
regarded as design labels until the resulting structures are independently
evaluated, for example by box counting or a radius-of-gyration method.

The numerical design was
\begin{equation}
\begin{aligned}
  &3\ \text{particle diameters}\times3\ \text{particle counts}\\
  &\quad\times3\ \text{requested fractal dimensions}
  \times5\ \text{random realizations}=135\ \text{cases}.
\end{aligned}
\end{equation}

\begin{center}
\begin{tabular}{|p{0.39\linewidth}|p{0.51\linewidth}|}
\hline
Design variable & Values \\
\hline
Primary-particle diameter $d_p$ & $1.0$, $1.5$, $2.0$ (simulation length units) \\
Number of particles $N$ & $50$, $100$, $200$ \\
Requested fractal dimension $D_f$ & $1.8$, $2.2$, $2.6$ \\
Random realizations & $5$ per parameter combination \\
Total number of cases & $135$ \\
\hline
\end{tabular}
\end{center}

Each realization was assigned a unique random seed. The generated particle
coordinates and radii were exported as CSV, LIGGGHTS data (\texttt{.sys}), and
VTK files. The LIGGGHTS data contain one granular atom type, no bonds, zero
initial translational and angular velocities, and an initial computational box
from $-100$ to $100$ in each coordinate direction.

\subsection{DEM relaxation and material parameters}

The generated structures were relaxed with LIGGGHTS using adhesive particle
contacts. Particle and contact data were stored at multiple time steps up to
step $100000$. The particle output contains positions, radii, particle IDs and
velocities, while the contact output contains the interacting particle IDs,
gap, overlap, contact area, contact state, normal and tangential forces, total
force and contact torque. These data were converted to VTK for inspection in
ParaView and subsequent graph construction.

The parameters confirmed directly from the generated particle data are listed
below. The remaining values must be copied from the final LIGGGHTS input before
this text is used in a report or publication.

\begin{center}
\begin{tabular}{|p{0.45\linewidth}|p{0.45\linewidth}|}
\hline
Property & Value used \\
\hline
Particle shape & Sphere \\
Particle diameter $d_p$ & $1.0$, $1.5$, $2.0$ (simulation length units) \\
Particle density $\rho_p$ & $2.0$ (in the selected LIGGGHTS unit system) \\
Initial translational/angular velocity & Zero \\
Bond types & None \\
Contact model & \texttt{[insert model from LIGGGHTS input]} \\
Young's modulus $E$ & \texttt{[insert verified value and unit]} \\
Poisson ratio $\nu$ & \texttt{[insert verified value]} \\
Coefficient of restitution $e$ & \texttt{[insert verified value]} \\
Sliding-friction coefficient $\mu$ & \texttt{[insert verified value]} \\
Rolling-friction coefficient $\mu_r$ & \texttt{[insert verified value]} \\
Hamaker constant $A_H$ & \texttt{[insert verified value in J]} \\
Minimum separation $h_{\min}$ & \texttt{[insert verified value and unit]} \\
Time step $\Delta t$ & \texttt{[insert verified value and unit]} \\
\hline
\end{tabular}
\end{center}

For each case, particle and contact VTK files were matched by time-step number.
A snapshot was classified as relaxed after three consecutive output intervals
satisfied both
\begin{equation}
 \frac{\Delta x_{\mathrm{RMS}}}{d_p}\leq 10^{-4},
 \qquad
 1-\frac{|C_k\cap C_{k+1}|}{|C_k\cup C_{k+1}|}\leq 0.01,
\end{equation}
where $C_k$ denotes the set of particle contacts at time step $k$. If no
snapshot fulfilled these criteria, the last common particle/contact time step
was retained and marked as non-converged.

\subsection{Graph construction and supervised targets}

Each relaxed agglomerate was represented as a graph. Primary particles form
the nodes and physical particle contacts form the edges. Every physical contact
was stored in both directions for message passing. Thus, an agglomerate with
$N$ particles and $N_c$ contacts has $N$ nodes and $2N_c$ directed edges.

The preprocessing stored three node features, but the minimal GNN used only
two structural node scalars:
\begin{equation}
 \boldsymbol{x}_i=
 \left[
   \frac{|\boldsymbol{x}_i-\boldsymbol{x}_c|}{R_g},\;
   z_i
 \right],
\end{equation}
where $d_{\mathrm{ref}}$ is twice the median particle radius,
$\boldsymbol{x}_c$ is the agglomerate center, $R_g$ is the radius of gyration,
and $z_i$ is the coordination number. Normalized, centered particle coordinates
were supplied separately to the encoder. The stored radius ratio was omitted
because it is approximately constant for the present monodisperse particles.

The converted DEM data provide 22 available edge/contact features, including
distance, overlap, gap, contact area, force and torque quantities and contact
flags. For the minimal model, only the normalized centre distance
$d_{ij}/d_{\mathrm{ref}}$ was selected as a stored edge scalar, because
connectivity and interparticle geometry are the quantities directly relevant
to the three purely structural targets considered here. Graph connectivity
and the normalized relative-position vector were additionally supplied for
message passing. Force, torque, overlap and contact-area fields were retained
in the processed data for sensitivity studies, but were not used by the
reported minimal model.

The convex hull of the particle centers was calculated with a small numerical
perturbation (Qhull option \texttt{QJ}). Hull vertices were classified as
surface particles. A breadth-first search through the contact graph then
defined the graph depth $h_i$ of every particle from this surface. Three
graph-level proxy targets were calculated:
\begin{align}
 \overline{h} &= \frac{1}{N}\sum_{i=1}^{N}h_i,\\
 f_{h\geq2} &= \frac{1}{N}\sum_{i=1}^{N}\mathbf{1}(h_i\geq2),\\
 S_{\mathrm{access}} &= \frac{1}{N}\sum_{i=1}^{N}\exp(-\beta h_i),
 \qquad \beta=0.5.
\end{align}
These quantities represent geometric shielding and accessibility proxies. They
are not yet oxygen-transport or oxygen-uptake results.

\subsection{GNN architecture and preliminary training}

The model embeds the two selected node scalars together with normalized
particle coordinates, and the selected edge-distance scalar together with
relative-position vectors, into a hidden space of dimension 64. Six residual
GINE message-passing layers with layer normalization and dropout are applied.
Global mean, maximum and size-normalized additive pooling produce one
aggregate-level representation, which is compressed to a latent vector
\begin{equation}
 \boldsymbol{\lambda}_A=[\lambda_1].
\end{equation}
Latent dimensions from one to four were tested. The selected dimension was the
smallest candidate whose mean out-of-fold $R^2$ was within 0.02 of the best
candidate. This rule selected one latent coordinate. A final two-layer
prediction head maps $\lambda_1$ to $\overline{h}$, $f_{h\geq2}$ and
$S_{\mathrm{access}}$; no precomputed aggregate-level descriptor is supplied
to the prediction head.

Training used Adam with a learning rate of $10^{-3}$, weight decay of $10^{-5}$,
at most 300 epochs, batch size 16, dropout 0.1 and random seed 7. Validation-based
early stopping and random rigid-rotation augmentation were used. Targets were
standardized using only the corresponding training data. All 135 graphs were
processed successfully. Five-fold cross-validation was grouped by parameter
combination, so that all random realizations of one $(d_p,N,D_f)$ combination
remained in the same held-out fold. The selected one-dimensional model gave:

\begin{center}
\begin{tabular}{|c|c|c|c|}
\hline
Fold & RMSE($\overline{h}$) & RMSE($f_{h\geq2}$) & RMSE($S_{\mathrm{access}}$) \\
\hline
1 & 0.315 & 0.051 & 0.044 \\
2 & 0.345 & 0.053 & 0.036 \\
3 & 0.452 & 0.054 & 0.041 \\
4 & 0.448 & 0.047 & 0.031 \\
5 & 0.393 & 0.040 & 0.033 \\
\hline
Mean & 0.390 & 0.049 & 0.037 \\
\hline
\end{tabular}
\end{center}

The combined out-of-fold $R^2$ values were 0.869 for mean graph depth, 0.894
for shielded fraction and 0.913 for accessibility score, with corresponding
MAE values of 0.293, 0.039 and 0.030. The mean $R^2$ across the three targets
was 0.892. Latent dimensions two, three and four gave mean $R^2$ values of
0.887, 0.888 and 0.876, respectively; hence additional latent coordinates did
not improve the result.

Particle-count-only baselines achieved mean out-of-fold $R^2$ values between
0.754 and 0.776, showing that the graph model captures structural variation
beyond particle count alone. A stricter baseline used particle count, radius
of gyration, contact count, contact density, mean coordination and maximum
coordination. Linear regression with these conventional descriptors achieved
a mean $R^2$ of 0.891, while a random forest achieved 0.909, slightly exceeding
the GNN by 0.016. Thus, the present proxy-target study demonstrates morphology
information beyond particle count, but it does not yet establish that detailed
graph encoding is necessary beyond well-chosen scalar descriptors. The more
decisive test will use resolved oxygen-transport targets for aggregates with
similar conventional descriptors but different internal topology.

The present results demonstrate the complete DEM-to-graph-to-GNN workflow, but
remain preliminary. The cross-validation latent values originate from
independently trained fold models and are therefore not interpreted as one
common latent coordinate. For descriptive latent analysis, final models are
trained on all graphs and their arbitrary sign and scale are aligned across
random seeds. Such all-data latent values are used only for interpretation,
not as leakage-free evidence of predictive performance.
